\section{Why NTT Matters}


Every lattice-based cryptographic operation involves polynomial multiplication in the ring
$R_q = \mathbb{Z}_q[X]/(X^n + 1)$. ML-KEM key generation performs $k^2$ polynomial
multiplications ($k = 2, 3,$ or $4$). ML-DSA signing performs roughly $k(\ell + k)$
multiplications per attempt, with multiple attempts due to rejection sampling. FHE
operations perform thousands of polynomial multiplications per bootstrapping.

Naive polynomial multiplication is $O(n^2)$: for $n = 256$, that is 65,536 multiply-and-add
operations, each with modular reduction. The NTT reduces this to $O(n \log n)$: 2,048
butterfly operations. For $n = 1024$ (FHE), the ratio is $1{,}048{,}576$ vs.\ $10{,}240$---a
factor of 100.

%% ============================================================
